\documentclass[12pt]{article}
\usepackage[utf8]{inputenc}
\usepackage{geometry}
\geometry{a4paper, margin=1in}

\title{B.Tech Syllabus, CSE\\
West Bengal University of Technology}
\date{2007-2011}

\begin{document}

\begin{titlepage}
\maketitle
\end{titlepage}

\tableofcontents

\section{M-101 [Credits 4]}
\subsection{Infinite Series}
Sequence, Convergence and Divergence of Infinite series – and typical examples of convergent and divergent series.\\ \\ 
Comparison test (statement only) and related problems.\\ \\
Ratio test (statement only) and related problems.\\ \\
Cauchy’s root test (statement only) and related problems.\\ \\
Alternating series, Leibnitz’s theorem (without proof), absolute convergence and related problems.
\subsection{Calculus of Functions of One Variable}
Review of limit and continuity and differentiability. \\ \\
Successive  differentiation,  Leibnitz’s  theorem  (without  proof  but  with  problems  of  the  type  of recurrence relations in derivatives of different orders and also to find $(y_n)_0)$.\\ \\
Rolle’s theorem (statement only); Mean Value Theorems—Lagrange \& Cauchy (statement only), Taylor’s  theorem  (without  proof  and  problems  in  respect  of  direct  use  and  applications  of  the theorem only), Expansions of functions by Taylor and Maclaurin series. Maclaurin’s expansion in infinite series of the functions: $\log (1+x)$, $e^x$ , $\sin x$, $\cos x$, $(a+x)^n$ , $n$ being a negative integer or a fraction L’Hospital’s Rule (statement only) and related problems.\\ \\
Integration of
\[
\int_0^{\pi/2}\cos^n(x)dx,\int_0^{\pi/2}\sin^n(x)dx,\int_0^{\pi/2}\cos^n(x)\sin^m(x)dx,\int_0^{\pi/2}\cos(nx)\sin(mx)dx
\]
$m,n$ are positive integers.
\subsection{Application}
Rectification
\subsection{Three Dimensional Geometry (Cartesian)}
Direction Cosine, Direction Ratio; Equation of a Plane (general form, normal form and intercept form); Equation of a Straight Line passing through one point and two points; Pair of intersecting planes representing a straight line.\\ \\
Elementary ideas of surfaces like sphere, Right Circular Cone and Right Circular Cylinder (through Geometrical configuration) and equations in standard forms.
\subsection{Calculus of Functions of Several Variables}
Introduction of Function of several variables and examples.\\ \\
Knowledge of limit and continuity. \\ \\
Partial derivative \& related problems. Homogeneous Functions and Euler’s Theorem (statement only) \& Problems upto 3 variables.\\ \\
Chain rules and related problems.\\ \\
Differentiation of implicit functions \& related problems.\\ \\
Total differentials and related problems.\\ \\
Maxima, minima and saddle points – definition, condition of extrema \& problems for two variables. Lagrange’s multiplier method – problems related to two variables only.\\ \\
Line Integral, Double Integrals, Triple Integral – Discussion w.r.t. different types of limits and problems; Moment of Inertia, Centre of Gravity.\\ \\
Jacobian – Definition and related problems for two variables.\\ \\
Applications to areas and volumes, surface area of revolution.
\subsection{Vector Calculus}
Scalar and Vector fields – Definition and Terminologies; Products: dot, cross, box, vector triple product.\\ \\
Gradient, directional derivative, divergence, curl (with problems).\\ \\
Tangent planes and normals and related problems.\\ \\
Statements of Green’s theorem, Divergence theorem, Stokes’ theorem with applications.

\section{M-201 [Credits 4]}
\subsection{Linear Algebra}
Introduction to the idea of a matrix; equality of matrices; special matrices. Algebraic operations of matrices: commutative property, associative property and distributive property. Transpose of a matrix (properties $(A^\top)^\top = A$, $(A+B)^\top = A^\top + B^\top$ , $(cA)^\top = cA^\top, (AB)^\top = B^\top A^\top$ to be stated (without proof) and verified by simple examples). Symmetric and Skew symmetric matrices.\\ \\
Properties of determinant (statement only); minor, co-factors and Laplace expansion of determinant; Cramer's rule and its application in solving system of linear equations of three variables.\\ \\
Singular and non-singular matrices; adjoint matrix; inverse of a matrix $(AB)^{-1} = B^{-1}A^{-1}$ to be stated and verified by example. Elementary row and column operations on matrices; definition of rank of a matrix; determination of rank of a matrix using definition.
\subsection{System of Linear Equations}
Consistency and Inconsistency. Gauss elimination process for solving a system of linear equations in three unknowns.
\subsection{Vector Space}
Basic idea of set, mapping, Binary Composition and Scalar field. Definition of vector space over the field of real numbers; Examples of vector space; Definition of sub-space of a vector space and a criterion for a sub-space; Definition of Linear combination, Linear independence and linear dependence of vectors with
examples. Definition of basis and dimension of vector space; Definition of Linear transformation: Definition of kernel and images of a Linear transformation; Kernel and Images of a Linear Transformation forming sub-spaces; Nullity and Rank of a Linear Transformation; Dim Ker T + Dim Im T = Dim V; Definition of Inner product space; Norm of a vector; Orthogonal and Orthonormal set of vectors.\\ \\
Eigenvalues and Eigenvectors of a matrix; Eigenvalues of a Real Symmetric Matrix; Necessary and
Sufficient Condition of diagonalization of matrices (statement only); Diagonalization of a matrix (problems restricted to $2\times 2$ matrix).
\subsection{Ordinary Differential Equations (ODE)}
Definition of order and degree of ODE; ODE of the first order: Exact equations; Definition and use of integrating factor; Linear equation and Bernoulli’s equation. ODE of first order and higher degree, simple problems.\\ \\ 
General ODE of 2nd order: D-operator method for finding particular integrals. Method of variation of parameters. Solution of Cauchy-Euler homogeneous linear equations. Solution of simple simultaneous linear differential equations.\\ \\
Verification of Legendre function ($P_n (x)$) and Bessel function ($J_n (x)$) as the solutions of Legendre and Bessel equations respectively. Graphical representations of these solutions.
\subsection{Laplace Transform (LT)}
Definition; Existence of LT; LT of elementary functions; First and second shifting properties; Change of scale property; LT of derivative of functions. LT of $(t^n f(t))$, LT of $f(t)/t^n$; LT of periodic function and unit step function. Convolution theorem (statement only).\\ \\
Inverse LT; Solution of ODE's (with constant coefficients) using LT.
\subsection{Numerical Methods}
Error: Absolute, Percentage, Relative errors. Truncation error, Round off error.\\ \\
Difference operator (forward, backward, central, shift and average operators); Different table, Propagation of Error. Definition of Interpolation and Extra-polation. Newton's forward and backward interpolation formula; Lagrange interpolation formula and corresponding error formulae (statement only).\\ \\
Numerical Differentiation: Using Newton’s forward and backward interpolation formula.\\ \\
Numerical Integration: Trapezoidal rule and Simpson's 1/3rd rule and corresponding error terms (statement only).
\section{M-301 [Credits 4]}
\subsection{Theory of Probability}
Random Experiment; Sample space; Random Events; Probability of events. Axiomatic   definition of probability; Frequency Definition of probability; Finite sample  spaces and equiprobable measure as special cases; Probability of Non-disjoint  events (Theorems). Counting techniques applied to probability problems; Conditional probability; General Multiplication Theorem; Independent events;  Bayes’ theorem and related problems.\\ \\
Random variables (discrete and continuous); Probability mass function; Probability density function and distribution function. Distributions: Binomial, Poisson, Uniform, Exponential, Normal, $t$ and $\chi^2$. Expectation and Variance ($t$ and $\chi^2$ excluded); Moment generating function; Reproductive Property of Binomal; Poisson and Normal Distribution (proof not required). Transformation of random variables (One variable); Chebychev inequality (statement) and problems.\\ \\
Binomial approximation to Poisson distribution and Binomial approximation to Normal distribution (statement only); Central Limit Theorem (statement); Law of large numbers (Weak law); Simple applications.
\subsection{Statistics}
Population; Sample; Statistic; Estimation of parameters (consistent and unbiased); Sampling distribution of sample mean and sample variance (proof not required).\\ \\
Point  estimate: Maximum likelihood estimate of statistical parameters (Binomial, Poisson and Normal distribution). Interval estimation.
\subsection{Testing of Hypothesis}
Simple  and  Composite  hypothesis; Critical Region; Level of Significance; Type I and Type II Errors; Best Critical  Region;  Neyman-Pearson  Theorem  (proof  not required);  Application  to  Normal  Population; Likelihood Ratio Test (proof not required); Comparison of Binomial Populations; Normal Populations; Testing of Equality of Means; $\chi^2$—Test of Goodness of Fit (application only).\\ \\
Simple idea of Bivariate distribution; Correlation and Regression; and simple problems. 

\section{M-401 [Credits 4]}
\subsection{Algebraic Structures}
Sets, Relation, Equivalence Relation, Equivalence Class \& Partition.\\ \\
Congruence Relation. Mapping; Inverse Mapping (Proof of Necessary and Sufficient Condition Excluded). Semigroup and Monoid; Group; Subgroup and Coset; Normal Subgroup; Quotient Group; Cycle Group, Permutation Group.\\ \\
Dihedral Group (upto $D_4$); Symmetric Group $S_3$, Homomorphism and Isomorphism; Modulo Group; Elementary Applications in Coding.\\ \\
Ring and Field: Ring; Subring; Morphism of Ring; Ideals and Quotient Ring. Integral Domain and Field; Finite Field; Statement of Relevant Theorems and  Examples.
\subsection{Lattice and Recurrence Relation}
Basic Idea; Sequence and Discrete function. Generating functions and application. 
\subsection{Advanced Graph Theory}
Planar and Dual Graphs. Kuratowski’s graphs. Homeomorphic graphs. Eulers formula $(n-e+r=2)$ for connected planar graph and its generalisation for graphs with connected components. Detection of planarity.\\ \\
Graph colouring. Chromatic numbers of $C_n$, $K_n$ , $K_{m,n}$ and other simple graphs. Simple applications of chromatic numbers. Upper bounds of chromatic numbers (Statements only). Chromatic polynomial. Statement of four and five
colour theorems.

\section{M(CS)-402 [Credits 4]}
\textbf{Operation Research and Optimization Techniques}
\subsection{Introduction}
Introduction to OR Modeling Approach and Various Real Life Situations.
\subsection{Linear Programming Problems (LPP)}
Basic LPP and Applications ; Various Components of LP Problem Formulation.
\subsection{Solving Linear Programming Problems}
Solving LPP: Using Simultaneous Equations and Graphical Method; Simplex Method; Duality Theory; Charnes’ Big–M Method. Transportation Problems and Assignment Problems.
\subsection{Network Analysis}
Shortest Path: Dijkstra Algorithm; Floyd Algorithm; Maximal Flow Problem (Ford-Fulkerson); PERT-CPM (Cost Analysis, Crashing, Resource Allocation excluded). 
\subsection{Inventory Control}
Introduction; EOQ Models; Deterministic and probabilistic Models; Safety Stock; Buffer Stock.
\subsection{Game Theory}
Introduction; 2-person Zero–sum Game; Saddle Point; Mini-Max and Maxi-Min Theorems (statement only); Games without saddle point; Graphical Method; Principle of Dominance.
\subsection{Queuing Theory}
Introduction; Basic Definitions and Notations; Axiomatic Derivation of the Arrival \& Departure (Poisson Queue). Pure Birth and Death Models; Poisson Queue Models: M/M/1 : $\infty$/FIFO and M/M/1: N/FIFO.
\end{document}